\section*{अध्याय \ref{ch:b2:climate-proxies} --- जलवायु अभिलेखागार और प्रॉक्सी}

\begin{solution}{exo:b2:climate-proxies:1}
(क) $\delta = (27.5 \times (-30) + 1350 \times 0)/(1350 + 27.5) = -825/1377.5 = -0.60$ प्रति मिल।
(ख) तब महासागर सारा जल रखेगा: $-0.60$ प्रति मिल, आज से \num{0.60} प्रति मिल
हल्का। (ग) $27.5/1377.5 = \qty{2.0}{\percent}$।
\end{solution}

\begin{solution}{exo:b2:climate-proxies:2}
(क) $T = 16.1 - 4.76\,(2.50 - (0 - 0.27)) = 16.1 - 4.76 \times 2.77 = 16.1 - 13.19 =
\qty{2.9}{\celsius}$। (ख) $T = 16.1 - 4.76\,(2.50 - 0.73) = 16.1 - 8.43 = \qty{7.7}{\celsius}$: यदि
जल भारी था तो वही सीपी लगभग पाँच अंश अधिक तापमान देती है, इसीलिए
जल का संघटन ज्ञात होना चाहिए। (ग) सीपी $2.71 - 2.50 = 0.21$ प्रति मिल भारी है,
इसलिए उसका जल $4.76 \times 0.21 = \qty{1.0}{\celsius}$ अधिक ठंडा था, \qty{1.9}{\celsius} पर।
\end{solution}

\begin{solution}{exo:b2:climate-proxies:3}
(क) $0.4\,\mathrm{e}^{0.9} = \qty{0.98}{mmol/mol}$ और $0.4\,\mathrm{e}^{1.8} =
\qty{2.42}{mmol/mol}$। (ख) गुणक $\mathrm{e}^{0.9} = 2.46$ है; प्रति अंश $\mathrm{e}^{0.09}
= 1.094$, यानी प्रति अंश \qty{9.4}{\percent} की वृद्धि। (ग) $T = \ln(\text{Mg/Ca}/0.4)/0.09$:
$\ln 5/0.09 = \qty{17.9}{\celsius}$ और $\ln 2.5/0.09 = \qty{10.2}{\celsius}$।
\end{solution}

\begin{solution}{exo:b2:climate-proxies:4}
(क) $2668 - 2213 + 1 = 456$ वर्ष। (ख) किसी एक वृक्ष में कोई वलय स्थानीय रूप से अनुपस्थित हो सकता है, या कोई ऐसा वलय
हो सकता है जो पढ़ा न जा सके; अनेक वृक्षों के बीच चौड़े और संकरे वलयों के प्रतिरूपों का मिलान
लुप्त वलयों को प्रकट करता है और हर वलय का वर्ष स्थिर करता है, और वृक्षों के बीच की सहमति
काल-निर्धारण की विश्वसनीयता को मापती है। (ग) वृक्ष-रेखा पर एक कारक, वृद्धि-ऋतु की ठंड,
सभी वृक्षों को सीमित करता है, इसलिए उनमें अच्छे और बुरे वर्षों का एक जैसा प्रतिरूप होता है; सौम्य, नम
निचले क्षेत्र में कोई एक कारक प्रभावी नहीं होता, वृक्षों के प्रतिरूप भिन्न होते हैं, और \omterm{def:b2:climate-proxies:dendro}{परस्पर काल-मिलान} विफल हो सकता है।
\end{solution}

\begin{solution}{exo:b2:climate-proxies:5}
(क) $125 \times 40/(3682.2 - 125) = 5000/3557.2 = 1.41$ प्रति मिल; $134 \times 40/3548.2 =
5360/3548.2 = 1.51$ प्रति मिल। (ख) $H - h = 3552.2$ के साथ: $3900/3552.2 = 1.10$, $4550/3552.2 = 1.28$
और $6500/3552.2 = 1.83$ प्रति मिल। (ग) $\delta_{\mathrm{ice}} = -1.0 \times 3552.2/130 = -27.3$ प्रति
मिल, किसी भी वर्तमान हिमचादर से भारी (ग्रीनलैंड $-30$ से $-35$, उत्तर अमेरिकी हिमचादर
$-40$, अंटार्कटिक बर्फ़ $-50$ तक कम)। यथार्थवादी बर्फ़ वाला द्रव्यमान संतुलन और छिद्र-जल
माप सहमत नहीं होते; लगभग आधे प्रति मिल का अंतर विधि की एक वास्तविक अनिश्चितता है।
\end{solution}

\begin{solution}{exo:b2:climate-proxies:6}
(क) $0.372\,\mathrm{e}^{0.0936T} = 0.4\,\mathrm{e}^{0.09T}$ से $\mathrm{e}^{0.0036T} =
0.4/0.372 = 1.0753$ मिलता है, अतः $T = \ln 1.0753/0.0036 = \qty{20.2}{\celsius}$। (ख)
\qty{2}{\celsius} पर: $0.4\,\mathrm{e}^{0.18} = 0.479$ और $0.372\,\mathrm{e}^{0.1872} = 0.449$, अतः
समंजित अनुपात $0.449/0.479 = 0.937$ है, \qty{6.3}{\percent} कम। (ग) $T = \ln(0.449/0.4)/0.09 =
\ln 1.1225/0.09 = \qty{1.3}{\celsius}$: पूर्वानुमानित नियम सीपी को लगभग
\qty{0.7}{\celsius} अधिक ठंडा पढ़ता है। ठंडे जल में अंशांकन का चुनाव एक अंश के स्तर पर
मायने रखता है।
\end{solution}

\begin{solution}{exo:b2:climate-proxies:7}
(क) होलोसीन से GI-1 तक, 2989 वर्षों के लिए \qty{112.19}{m} बर्फ़: $11\,219/2989 =
\qty{3.75}{cm}$ प्रति वर्ष; GI-16 से GI-17 तक, 1160 वर्षों के लिए \qty{17.89}{m}: $1789/1160 =
\qty{1.54}{cm}$ प्रति वर्ष। (ख) बर्फ़
बहती है: जैसे-जैसे हर परत दबती है, वह पार्श्व में फैलती और पतली होती है, और गहरी, पुरानी परतें
अधिक समय तक पतली होती रही हैं। (ग) जब परतें बहुत पतली हो जाती हैं, एक सेंटीमीटर या उससे कम, तो उपकरण
एक वर्ष को अगले से अलग नहीं कर पाते, और \omterm{def:b2:climate-proxies:age-model}{आयु-मॉडल} को उसके बजाय बंधन बिंदुओं या प्रवाह मॉडलों का
उपयोग करना पड़ता है।
\end{solution}

\begin{solution}{exo:b2:climate-proxies:8}
(क) अंतराल का अंश $(1550 - 1492.45)/(1604.64 - 1492.45) = 57.55/112.19 = 0.513$ है,
अतः आयु $11\,703 + 0.513 \times 2989 = 11\,703 + 1533 = \num{13236}$ वर्ष b2k है। (ख) 1950, 2000 से 50
वर्ष पहले है: 1950 से \num{13186} वर्ष पहले। (ग) $50 + 0.513 \times (93 - 50) = 72$ वर्ष।
\end{solution}

\begin{solution}{exo:b2:climate-proxies:9}
(क) हर अनिश्चित परत आधा वर्ष योगदान करती है: $2 \times 93 = 186$ और $2 \times 1287 = 2574$
अनिश्चित परतें। (ख) $\frac12\sqrt{186} = 6.8$ वर्ष और $\frac12\sqrt{2574} = 25.4$ वर्ष।
(ग) $93/6.8 = 14$ और $1287/25.4 = 51$। $\frac12 n$ और $\frac12\sqrt{n}$ का अनुपात
$\sqrt{n}$ है: वह अनिश्चित परतों की संख्या के वर्गमूल के रूप में बढ़ता है, इसलिए दोनों नियमों के बीच
चुनाव गहराई के साथ अधिकाधिक मायने रखता है।
\end{solution}

\begin{solution}{exo:b2:climate-proxies:10}
(क) \omterm{prop:b2:isotope-geochemistry:mwl}{उल्कीय जल रेखा} पर $\delta$D, $\delta^{18}$O से आठ गुना बदलता है, इसलिए
प्रवणता $8 \times 0.8 = 6.4$ प्रति मिल $\delta$D प्रति अंश है। गिरावट $441.7 - 396.6 = 45.1$
प्रति मिल है, अतः शीतलन $45.1/6.4 = \qty{7.0}{\celsius}$। (ख) हिमनदीय महासागर लगभग
$8 \times 1.0 = 8$ प्रति मिल भारी था, $\delta$D में; आज के महासागर से बनी हिम
उतनी ही हल्की होती, $-441.7 - 8 = -449.7$ प्रति मिल। गिरावट $53.1$ प्रति मिल हो जाती है और शीतलन
$53.1/6.4 = \qty{8.3}{\celsius}$। (ग) $9.39 - 0.09 = \qty{9.30}{\celsius}$, अतः $45.1/9.30 = 4.85$
प्रति मिल प्रति अंश। प्रकाशित श्रेणी में हिमचादर की बदलती ऊँचाई के लिए भी सुधार किया गया है,
और उसका तापमान-संबंध जल-समस्थानिकों का अनुकरण करने वाले मॉडलों से आता है, केवल
स्थानिक प्रवणता से नहीं।
\end{solution}

\begin{solution}{exo:b2:climate-proxies:11}
(क) अनुपातों के लघुगणक $-0.528$, $-0.223$, $0.191$, $0.501$, $0.904$, $1.227$ हैं; उनका
माध्य $0.345$ है और माध्य तापमान 14। $S_{TT} = \sum (T - 14)^2 = 280$ और $S_{T\ell} =
\sum (T - 14)(\ell - 0.345) = 24.93$ के साथ, $A = 24.93/280 = \qty{0.0890}{\per\celsius}$ और $\ln B =
0.345 - 0.0890 \times 14 = -0.901$, अतः $B = \qty{0.406}{mmol/mol}$। (ख) अवशेष
$+0.017$, $-0.034$, $+0.023$, $-0.023$, $+0.025$ और $-0.009$ हैं, अनुपात के कुछ प्रतिशत, बिना किसी
प्रवृत्ति के; समंजन $0.4\,\mathrm{e}^{0.09T}$ के निकट है। (ग) $T = (\ln 1.30 + 0.901)/0.0890 = (0.262 +
0.901)/0.0890 = \qty{13.1}{\celsius}$। (घ) $T = (\ln 0.45 + 0.901)/0.0890 = (-0.799 +
0.901)/0.0890 = \qty{1.2}{\celsius}$। यह मान अंशांकन के सबसे ठंडे नमूने,
\qty{4}{\celsius}, से नीचे है: नियम का बहिर्वेशन हुआ है, और निम्न तापमानों पर सीपियों का Mg/Ca
कार्बोनेट आयन के प्रभाव में घातांकी नियम से हटने के लिए जाना जाता है।
\end{solution}

\begin{solution}{exo:b2:climate-proxies:12}
(क) महासागर के प्रति इकाई क्षेत्रफल, वर्तमान महासागर का जल, $H$ जो 0 पर है, भंडारों
$h_i$, जो $\delta_i$ पर हैं, और शेष महासागर $H - h$ के बीच बँटता है, जहाँ $h = \sum h_i$: $0 = \sum h_i\delta_i +
(H - h)\,\Delta\delta$। पहले चरण के बाद $\Delta\delta = 60 \times 30/(3682.2 - 60) =
1800/3622.2 = 0.50$ प्रति मिल; दूसरे के बाद $(1800 + 70 \times 45)/(3682.2 - 130) =
4950/3552.2 = 1.39$ प्रति मिल। (ख) $-4950/130 = -38.1$ प्रति मिल। (ग) नियम $1.39/130 = 0.0107$
प्रति मिल प्रति मीटर है, और पहले चरण के लिए $0.50/0.0107 = \qty{47}{m}$ देता है, \qty{60}{m} के
बजाय: पहली बर्फ़ माध्य से भारी थी, इसलिए उसके हर मीटर ने महासागर को
कम खिसकाया। एक अवस्था पर अंशांकित रैखिक नियम दूसरी अवस्थाओं को गलत पढ़ता है जब
बर्फ़ का माध्य संघटन बदलता है।
\end{solution}

\begin{solution}{pb:b2:climate-proxies:1}
\textbf{1.} $361.84 \times 10^6\ \unit{km^2} \times \qty{0.130}{km} = \qty{4.70e7}{km^3}$, लगभग 47
मिलियन घन किलोमीटर।
\textbf{2.} $130 \times 40/(3682.2 - 130) = 5200/3552.2 = 1.46$ प्रति मिल।
\textbf{3.} $5000/3557.2 = 1.41$ और $5360/3548.2 = 1.51$ प्रति मिल।
\textbf{4.} $4550/3552.2 = 1.28$ और $6500/3552.2 = 1.83$ प्रति मिल। बर्फ़ का संघटन
उत्तर को \num{0.55} प्रति मिल खिसकाता है, समुद्र-तल \num{0.10}: बर्फ़ अधिक मायने रखती है।
\textbf{5.} $\delta_{\mathrm{ice}} = -1.0 \times 3552.2/130 = -27.3$ प्रति मिल।
\textbf{6.} $40/3682.2 = 0.0109$ प्रति मिल प्रति मीटर, जो दिनांकित समुद्र-तलों के \num{0.011} तक
पूर्णांकित होता है।
\textbf{7.} स्थिर $\delta^{18}\mathrm{O}_{\mathrm{sw}}$ पर, $\mathrm{d}T = -4.76\,\mathrm{d}\delta^{18}
\mathrm{O}_{\mathrm{b}}$, अतः $\mathrm{d}T = -1$ से $\mathrm{d}\delta^{18}\mathrm{O}_{\mathrm{b}} =
1/4.76 = 0.21$ प्रति मिल मिलता है।
\textbf{8.} $\delta^{18}\mathrm{O}_{\mathrm{b}} = (16.1 - T)/4.76 + \delta^{18}\mathrm{O}_{\mathrm{sw}}
- 0.27 = 14.1/4.76 - 0.27 = 2.96 - 0.27 = 2.69$ प्रति मिल।
\textbf{9.} $17.1/4.76 + 1.0 - 0.27 = 3.59 + 0.73 = 4.32$ प्रति मिल। परिवर्तन $4.32 - 2.69 =
1.63$ प्रति मिल है: \num{1.00} जल से और $3.0 \times 0.21 = 0.63$
\qty{3}{\celsius} के शीतलन से।
\textbf{10.} $T = \ln(2.0/0.4)/0.09 = \ln 5/0.09 = \qty{17.9}{\celsius}$।
\textbf{11.} $0.18/0.1 = \qty{1.8}{\celsius}$ का शीतलन।
\textbf{12.} Mg/Ca कैल्साइट में दो तत्वों का अनुपात है, जो कैल्सीकरण के तापमान से
तय होता है; उसमें ऑक्सीजन के समस्थानिक शामिल नहीं हैं, और इसलिए वह जल के उस संघटन की
उपेक्षा करता है जिसे बर्फ़ बदलती है।
\textbf{13.} $4.92 - 3.23 = 1.69$ प्रति मिल।
\textbf{14.} $\Delta = h \times 40/(H - h)$ को $h$ के लिए हल करने पर: $h = H\Delta/(\Delta + 40) = 3682.2
\times 1.69/41.69 = \qty{149}{m}$, जो समुद्र-तल के साक्ष्य के 125 से \qty{134}{m} से अधिक है:
परिवर्तन का कुछ भाग शीतलन होना चाहिए।
\textbf{15.} $1.69 - 1.46 = 0.23$ प्रति मिल, और $4.76 \times 0.23 = \qty{1.1}{\celsius}$।
\textbf{16.} $1.69 - 1.00 = 0.69$ प्रति मिल, और $4.76 \times 0.69 = \qty{3.3}{\celsius}$ (3.28)।
\textbf{17.} $0.23/0.25 = \qty{0.9}{\celsius}$ और $0.69/0.25 = \qty{2.8}{\celsius}$।
\textbf{18.} उत्तर अटलांटिक गहरा जल 2--\qty{4}{\celsius} से \qty{-1.1}{\celsius} तक ठंडा हुआ,
3.1 से \qty{5.1}{\celsius}; अंटार्कटिक तली जल 0--\qty{1}{\celsius} से
\qty{-1.9}{\celsius} तक, 1.9 से \qty{2.9}{\celsius}; उत्तर-पश्चिम अटलांटिक का Mg/Ca
\qty{1.8}{\celsius} देता है। सभी लगभग 2 और \qty{5}{\celsius} के बीच हैं, छिद्र-जल सुधार के 2.8 से
\qty{3.3}{\celsius} के आसपास, और द्रव्यमान संतुलन के \qty{1}{\celsius} से ऊपर, जिसमें बर्फ़ $-40$ प्रति
मिल पर है।
\textbf{19.} $-396.6 - (-441.7) = 45.1$ प्रति मिल।
\textbf{20.} $8 \times 0.8 = 6.4$ प्रति मिल प्रति अंश, अतः $45.1/6.4 = \qty{7.0}{\celsius}$।
\textbf{21.} $9.39 - 0.09 = \qty{9.30}{\celsius}$, और $45.1/9.30 = 4.85$ प्रति मिल प्रति अंश।
\textbf{22.} डोम~C, \qty{9.3}{\celsius}, वैश्विक सतह, 5 से \qty{7}{\celsius}, से अधिक, जो
गहरे महासागर, लगभग \qty{3}{\celsius}, से अधिक है।
\textbf{23.} गहरा जल आज भी ठंडा है, 0 से \qty{4}{\celsius}, और समुद्री जल अपने हिमांक से बहुत
नीचे ठंडा नहीं हो सकता, जिस तक हिमनदीय तली जल \qty{-1.9}{\celsius} पर लगभग पहुँच गया था:
गहरे महासागर के पास खोने को केवल कुछ अंश थे, जबकि अंटार्कटिका के ऊपर की वायु की
ऐसी कोई निचली सीमा नहीं है।
\textbf{24.} छिद्र-जल सुधार के साथ, अंतिम हिमनदीय अधिकतम पर गहरा महासागर आज की तुलना में लगभग \textbf{\qty{3.3}{\celsius}
अधिक ठंडा} था, स्वतंत्र गहरे-जल अनुमानों के 2 से \qty{5}{\celsius} के
निकट।
\end{solution}
